\section{Introduction}\label{ap:introduction}

Erd\H{o}s asked whether every set $A$ of positive integers for which
$\sum_{a\in A}1/a$ diverges contains arithmetic progressions of every
finite length~\cite[Problem~4.33.6]{Erdos1974}.
We prove a quantitative density bound that implies this assertion.
For integers $k\ge3$ and $N\ge1$, let $r_k(N)$ be the largest
cardinality of a subset of $[N]=\{1,\ldots,N\}$ containing no
nonconstant $k$-term arithmetic progression
$a,a+d,\ldots,a+(k-1)d$, with $d>0$.
All logarithms are natural.

\begin{theorem}\label{ap:main}
For each fixed integer $k\ge3$ there are constants
$C_k,c_k,\varepsilon_k>0$ such that, for every $N\ge2$,
\begin{equation}\label{ap:density-bound}
 r_k(N)\le C_kN\exp\bigl(-c_k(\log N)^{\varepsilon_k}\bigr).
\end{equation}
Equivalently, there is $A_k\ge1$ such that an $\alpha$-dense subset
of $[N]$ contains a nonconstant $k$-term progression whenever
\begin{equation}\label{ap:threshold}
 \log N\ge A_k\bigl(2+\log(1/\alpha)\bigr)^{A_k},
 \qquad 0<\alpha\le1.
\end{equation}
\end{theorem}


The constants may depend on the fixed length $k$; the exponent is
not optimized. The reciprocal-sum consequence follows immediately
from summing the density bound over dyadic intervals.

\begin{corollary}[Divergent reciprocal sums]\label{ap:reciprocal}
Every set $A\subseteq\mathbb N$ with
$\sum_{a\in A}a^{-1}=\infty$ contains nonconstant arithmetic
progressions of every finite length.
\end{corollary}

\begin{proof}
Fix $k\ge3$, and suppose that $A$ contains no nonconstant
$k$-term progression. For $m\ge0$, translation identifies
$[2^m,2^{m+1})\cap\mathbb Z$ with an interval of length $2^m$.
Thus $|A\cap[2^m,2^{m+1})|\le r_k(2^m)$. Theorem~\ref{ap:main}
implies, for $m\ge1$,
\[
 \sum_{a\in A\cap[2^m,2^{m+1})}\frac1a
 \le 2^{-m}r_k(2^m)
 \le C_k\exp\bigl(-c_k(m\log2)^{\varepsilon_k}\bigr).
\]
For every $d,\varepsilon>0$, the sequence $e^{-dm^\varepsilon}$
is summable: since $m^\varepsilon/\log m\to\infty$, its terms
are at most $m^{-2}$ for all sufficiently large $m$.
The $m=0$ interval contributes at most one. The reciprocal sum of
$A$ is therefore finite, a contradiction. Divergence also implies
that $A$ is infinite, which supplies lengths one and two.
\end{proof}

\subsection{Density thresholds and summability}

The threshold in \eqref{ap:threshold} is quasipolynomial in
$1/\alpha$: the logarithm of a sufficient interval length is a fixed
